% TOP_PROBLEM: {"id": 30006846, "title": "Woodin's HOD Conjecture", "category_id": 10, "category": {"id": 10, "name": "set_theory", "display_name": "Set Theory", "description": "Foundations of mathematics, infinite sets, and cardinality.", "slug": "set-theory", "order_index": 10, "created_at": "2026-07-31T15:26:25.671Z"}, "status": "open"}
% FIRST_POSTED: 2026-09-27
% !TeX program = lualatex
% ATTEMPT_STATUS: unresolved
\documentclass[11pt]{article}
\usepackage[margin=25mm]{geometry}
\usepackage{fontspec}
\setmainfont{Latin Modern Roman}
\usepackage{amsmath,amssymb,mathtools}
\usepackage{unicode-math}
\setmathfont{Latin Modern Math}
\usepackage{xurl,hyperref}
\hypersetup{colorlinks=true,urlcolor=blue}
\setcounter{secnumdepth}{0}
\setlength{\parindent}{0pt}
\setlength{\parskip}{5pt}
\setlength{\emergencystretch}{3em}
\begin{document}
\section*{\texorpdfstring{Woodin's HOD Conjecture}{Woodin's HOD Conjecture}}\label{30006846}
\subsection{Definitions and mathematical statement}
Work in ZFC. The inner model $\mathrm{HOD}$ consists of sets hereditarily definable from ordinal parameters. For an uncountable regular cardinal $\lambda$, put $S^\lambda_\omega=\{\alpha<\lambda:\operatorname{cf}^V(\alpha)=\omega\}$. A set is stationary in $\lambda$ if it meets every closed unbounded subset of $\lambda$. Say $\lambda$ is \emph{$\omega$-strongly measurable in HOD} if some $\kappa<\lambda$ satisfies $(2^\kappa)^{\mathrm{HOD}}<\lambda$ and no partition of $S^\lambda_\omega$ into $\kappa$ sets, represented by a sequence in HOD, has every part stationary in $V$. A cardinal $\delta$ is extendible if for every $\alpha>\delta$ there are $\beta$ and an elementary $j:V_\alpha\to V_\beta$ with critical point $\delta$ and $j(\delta)>\alpha$. The selected conjecture is that existence of an extendible cardinal implies
\[
 \forall\gamma\ \exists\lambda>\gamma\quad
 [\lambda\text{ regular and not $\omega$-strongly measurable in HOD}].
\]
Thus the stationary-set test is external, in $V$, while the partition sequence must belong to HOD.
\par\smallskip\textit{Formulation and concept references:} \href{http://panhellenic-logic-symposium.org/12/slides/Day2_Bagaria.pdf}{[S1]}, \href{https://arxiv.org/html/1605.00613v1}{[E1]}.

\subsection{Short English statement}
Assuming an extendible cardinal exists, must there be arbitrarily large regular cardinals at which HOD fails the specified strong measurability property?
\subsection{Research attempt: external stationarity and the partition quantifiers}
\textbf{Status and a literal convention issue.} The intended HOD conjecture is not proved here. There is a small but consequential convention in the recorded definition: the witness $\kappa$ must be nonzero, as in the usual use of an infinite cardinal in this definition. If the bare inequality $\kappa<\lambda$ is read as allowing $\kappa=0$, then $(2^0)^{\mathrm{HOD}}=1<\lambda$ and the nonempty set $S^\lambda_\omega$ has no partition into zero pieces. That would make every uncountable regular $\lambda$ satisfy the displayed strong-measurability definition for a vacuous reason, destroying the intended conjecture. The inherited statement is preserved above; this attempt explicitly uses nonzero witnesses, in particular the standard infinite-cardinal convention, and does not present that wording defect as a solution of the intended problem.

The primary source [E1], Definition 9 and the subsequent conjecture, confirms that stationarity is tested in $V$ while the partition sequence belongs to HOD. The source check on 27 September 2026 also examined the primary compactness results of Goldberg and Poveda [E2]. Those have additional hypotheses about stationary sets of smaller cardinals and do not themselves give the proper-class conclusion selected here.

\textbf{Elementary facts that fix the universe of each assertion.} Write $M=\mathrm{HOD}^V$. If $\lambda$ is a regular cardinal in $V$, it is a regular cardinal in $M$. Indeed, an $M$-bijection from a smaller ordinal onto $\lambda$, or an $M$-cofinal map from a smaller ordinal into $\lambda$, would also be such a map in $V$, contradicting cardinality or regularity there. The reverse implication is not supplied by this argument: $V$ may have functions absent from $M$.

The set
\[
 S=S^\lambda_\omega=\{\alpha<\lambda:\operatorname{cf}^V(\alpha)=\omega\}
\]
itself belongs to HOD. It is definable in $V$ from the ordinal $\lambda$, and all its members, and their transitive closures, are ordinals. Nevertheless $M$'s internally computed set of ordinals of cofinality $\omega$ need not be this same set. A countable cofinal sequence in $V$ need not belong to $M$. Thus the fact that $S\in M$ does not allow one to replace its externally specified definition by an internal cofinality calculation.

For an uncountable regular $\lambda$, $S$ is stationary in $V$. Given a club $C\subseteq\lambda$, choose a strictly increasing sequence $\alpha_n\in C$. Regularity gives $\alpha=\sup_n\alpha_n<\lambda$, closure gives $\alpha\in C$, and strict increase gives $\operatorname{cf}^V(\alpha)=\omega$. Thus $\alpha\in C\cap S$.

If $C\in M$ is a club in $\lambda$ according to $M$, it is a club according to $V$ as well. Unboundedness is a statement about the same ordinals below $\lambda$. Closure can be expressed by the condition that every ordinal $\beta<\lambda$ with $\sup(C\cap\beta)=\beta$ belongs to $C$, again using the same set and ordinals. It follows that
\[
 A\in M\text{ stationary in }V
 \quad\Longrightarrow\quad M\models\text{$A$ is stationary in $\lambda$}.
 \tag{303.1}
\]
The converse cannot be inferred: internal stationarity tests only clubs belonging to $M$, whereas external stationarity tests all clubs in $V$. No assumption in the selected statement identifies those two collections of clubs.

\textbf{The exact negation at one regular cardinal.} Under the nonzero convention, failure of strong measurability at $\lambda$ means
\[
 \begin{split}
 \forall\kappa\in(0,\lambda)\quad
 (2^\kappa)^M<\lambda\ \Longrightarrow\ 
 \exists\langle S_\xi:\xi<\kappa\rangle\in M\quad
 &S=\bigsqcup_{\xi<\kappa}S_\xi,\\
 &\text{every }S_\xi\text{ is stationary in }V.
 \end{split}
 \tag{303.2}
\]
If witnesses are restricted to infinite cardinals, the same restriction is made in this formula. Producing just one partition into two stationary pieces is not (303.2). It eliminates one possible witness, but there may be another eligible $\kappa$ for which the required partition is absent.

The quantifier requiring the entire sequence to lie in $M$ is also stronger than requiring each individual part to lie in $M$. In general a sequence of HOD objects need not be in HOD. For example, if $V$ has a real $r$ outside HOD, the sequence of its zero-one digits has every entry in HOD, but the sequence itself is not in HOD, since it reconstructs $r$. This conditional example isolates the logical distinction without assuming that such a real exists in every model of ZFC.

\textbf{Lemma 303.1 (coarsening is legitimate).} Suppose $0<\kappa\leq\mu\leq\lambda$ and an $M$-sequence partitions $S$ into $\mu$ sets which are stationary in $V$. Then there is an $M$-sequence partitioning $S$ into $\kappa$ sets stationary in $V$.

\emph{Proof.} Use the ordinal-definable surjection $f:\mu\to\kappa$ given by $f(\xi)=\xi$ for $\xi<\kappa$ and $f(\xi)=0$ otherwise. If the original parts are $S_\xi$, let
\[
 T_\eta=\bigcup\{S_\xi:f(\xi)=\eta\}\qquad(\eta<\kappa).
\]
This sequence is definable from the original sequence and ordinals, so belongs to $M$. Its parts are disjoint and cover $S$. Each contains $S_\eta$, which is stationary in $V$, hence is itself stationary there. $\square$

In particular an $M$-partition into $\lambda$ externally stationary sets would prove that $\lambda$ is not strongly measurable, independently of the exponentiation condition. Failure of a partition at some $\kappa$ rules out all finer partitions with at least $\kappa$ pieces; the converse implication from the absence of a finer partition to the absence of a coarser one is invalid.

\textbf{A complete conditional case.} If $V=\mathrm{HOD}$, then every uncountable regular $\lambda$ fails the intended strong-measurability property. Use Solovay's stationary splitting theorem: every stationary subset of an uncountable regular $\lambda$ can be partitioned into $\lambda$ stationary sets. This standard theorem is explicitly recorded in [E1, Theorem 13] and is used here as an external input. Apply it to $S$. When $V=M$, the resulting sequence belongs to $M$ and its stationarity is exactly the required external stationarity. Lemma 303.1 supplies every eligible smaller partition, proving (303.2).

Thus under the additional assumption $V=\mathrm{HOD}$, the proper-class conclusion follows without needing extendibility. This is a conditional subcase, not an argument that extendibility forces $V=\mathrm{HOD}$, nor a claim that ZFC proves existence of a model with all desired large cardinals and that extra equality.

\textbf{The club filter explains the measurability obstruction.} The following conditional calculation is useful independently of any dichotomy theorem. Suppose $S_0\in M$ is stationary in $\lambda$ in $V$ and has no partition into two $M$-sets both stationary in $V$. Define
\[
 U=\{A\in M:A\subseteq S_0\text{ and some club }C\subseteq\lambda
                    \text{ satisfies }C\cap S_0\subseteq A\}.
 \tag{303.3}
\]
The club in this definition is allowed to be in $V$. For every $A\in M$ contained in $S_0$, at least one of $A$ and $S_0\setminus A$ is nonstationary in $V$, by the assumed lack of a split. Its complement in $S_0$ therefore contains a club trace. Both sides cannot be nonstationary, since the intersection of two clubs would then avoid $S_0$. Thus $U$ decides exactly one side of each such complementary pair.

Intersections of fewer than $\lambda$ clubs are club. To see unboundedness for a family of size $\theta<\lambda$, repeatedly choose a new ordinal above the current one and above a point from each club; at the supremum of countably many such steps, each club is met cofinally below the same limit. Regularity keeps all suprema below $\lambda$, and closure puts the limit in every club. Closure of the intersection is immediate. It follows that $U$ is $\lambda$-complete for sequences in $M$. It contains no singleton because a tail club avoids any given point.

Finally, $U\in M$: membership in (303.3) is definable in $V$ from $S_0$ and $\lambda$; since $S_0$ is ordinal definable, $U$ is ordinal definable, and each of its elements is hereditarily ordinal definable. Stationarity makes $S_0$ have size $\lambda$ in $V$, and it also has that size in $M$, since any smaller $M$-enumeration would be smaller in $V$. An $M$-bijection between $\lambda$ and $S_0$ transports $U$ to a nonprincipal $\lambda$-complete ultrafilter on $\lambda$ inside $M$.

This calculation explains why an unsplittable externally stationary set can yield internal measurability. It does not construct such a set from the selected hypotheses, and conversely an arbitrary HOD measure is not automatically an obstruction to all required external stationary partitions.

\textbf{Where the extendibility route stops.} An extendibility embedding has the form $j:V_\alpha\to V_\beta$. It is not, merely by its definition, a global elementary embedding of HOD into itself. Ordinal definitions can quantify over ranks outside the chosen segment, and the image of a sequence must still be checked against clubs in the relevant ambient universe. A proof using $j$ to obtain or rule out a stationary partition must control these domain, definability, and stationarity issues explicitly. Invoking the HOD dichotomy also does not select the desired side of that dichotomy without an additional argument.

\textbf{Checks and outcome.} The accompanying diagnostic checks coarsening of finite partitions and records the directions of the quantifiers in (303.1)--(303.2). It is a check of finite set identities only; finite sets have no analogue that verifies the uncountable stationarity claims used here. Those claims are handled by the proofs and explicitly named splitting theorem above.

The attempt proves the stated conditional $V=\mathrm{HOD}$ case, coarsening reductions, and the filter consequence of an unsplittable HOD set, while exposing the zero-witness convention and the external/internal distinction. The unresolved step is deriving arbitrarily large regular cardinals satisfying all of (303.2) from extendibility alone. None of the elementary arguments or restricted compactness results cited here supplies that implication. The intended conjecture remains unresolved by this attempt.

\subsection{Sources}
\begin{itemize}
\item[S1]Formal statement, Slide 'The HOD Conjecture' (invited talk, 12th Panhellenic Logic Symposium, 2019).
\url{http://panhellenic-logic-symposium.org/12/slides/Day2_Bagaria.pdf}.
\textit{Formulation source.}
\item[E1]W. Hugh Woodin, Jacob Davis, and Daniel Rodríguez, The HOD Dichotomy, arXiv:1605.00613v1 (2016), Sections 1 and 3.
\url{https://arxiv.org/html/1605.00613v1}.
\textit{Added primary source: Extendible cardinals, omega-strong measurability in HOD, and the proper-class formulation of the HOD conjecture.. Consulted 24 September 2026.}
\item[E2]Gabriel Goldberg and Alejandro Poveda, \emph{Compactness phenomena in HOD}, author-hosted primary manuscript.
\url{https://math.berkeley.edu/~goldberg/Papers/CompactnessPhenomenaInHOD.pdf}.
\textit{Consulted 27 September 2026; the conditional compactness results about singular cardinals are not treated as a resolution of the proper-class HOD conjecture.}
\end{itemize}
\end{document}
